\documentclass[11pt]{article}
\usepackage{amsmath,amsthm,booktabs,longtable,array,enumitem,needspace}
\usepackage{hyperref}
\usepackage{komorebi-publication}
\renewcommand{\KomorebiTitle}{Conductors, Native Markings, and the Equal-Capacity Boundary}
\renewcommand{\KomorebiSubtitle}{Toward Arithmetic Packet Realization\par Part III of the Conjecture 66 progress paper}
\renewcommand{\KomorebiSeries}{Paper XI / Research companion}
\renewcommand{\KomorebiKind}{Publication Block 04}
\renewcommand{\KomorebiEdition}{Working continuation v0.4.0 / 10 October 2026}
\renewcommand{\KomorebiShortTitle}{Paper XI / Conjecture 66 / Part III}
\renewcommand{\KomorebiTitleSize}{28}
\renewcommand{\KomorebiTitleLeading}{33}
\KomorebiSetup
\newtheorem{proposition}{Proposition}[section]
\newtheorem{lemma}[proposition]{Lemma}
\newtheorem{corollary}[proposition]{Corollary}
\theoremstyle{definition}\newtheorem{definition}[proposition]{Definition}
\newcommand{\GL}{\operatorname{GL}}
\newcommand{\End}{\operatorname{End}}
\newcommand{\Ad}{\operatorname{Ad}}
\newcommand{\Hom}{\operatorname{Hom}}
\newcommand{\Sw}{\operatorname{Sw}}
\newcommand{\Arf}{\operatorname{Arf}}
\newcommand{\Tr}{\operatorname{Tr}}
\newcommand{\AI}{\operatorname{AI}}
\newcommand{\Ind}{\operatorname{Ind}}
\newcommand{\F}{\mathbf F}
\newcommand{\C}{\mathbf C}
\newcommand{\Z}{\mathbf Z}
\newcommand{\OO}{\mathcal O}
\newcommand{\pp}{\mathfrak p}
\newcommand{\A}{\mathfrak A}
\newcommand{\Pp}{\mathfrak P}
\newcommand{\source}[1]{\par{\small\textsf{Source and status.}\ #1}\par}
\setcounter{section}{21}
\begin{document}
\KomorebiCover
\section*{Reading this continuation}
\addcontentsline{toc}{section}{Reading this continuation}
A conductor is an arithmetic invariant, not a complete description of a type. A nondegenerate commutator is a finite invariant, not a complete description of a marked dyadic extension. These two distinctions meet at the smallest dyadic construction with exactly enough quotient capacity. There is no spare plane there to absorb an independently prescribed quadratic mismatch.

This continuation audits the conductor and marking branches at their actual hypotheses. It keeps three arithmetic objects separate: the original source representation, a candidate obtained through a classical induction and descent construction, and a candidate arising from a proposed explicit type. A comparison is useful only when its normalization and its subject are fixed.

The resulting finite theory is also constructive. A discrepancy between two quadratic refinements becomes one linear functional. When its Arf obstruction vanishes, a single transvection gives an explicit correction. For a supplied odd-order finite action, a further central discrepancy can be removed by a finite sum. Neither correction constructs the missing arithmetic action, chooses a canonical native marking, or settles the full phase equations.

\textbf{Publication status.} This is an independently editable Part III. The first two parts retain their prior source and review boundaries. In particular, the 33-page Block 02 reader remains a preserved PDF because its complete editable archive was not recovered. The original three Paper XI readers have not been revised here. The current ledger records mathematical obligations separately from documentary gaps.

\textbf{Evidence status.} Ordinary finite proofs are supplied below. The arithmetic construction records are inherited arguments audited against selected primary statements. New exact tests check finite models and normalizations. No fresh Lean compilation or external peer review is claimed.

\tableofcontents
\clearpage
\section{Conductor conventions and actual order periods}\label{sec:conventions}
\subsection{Four numbers, not one unnamed conductor}
Let $F$ be a nonarchimedean local field with residue cardinality $q$, and let $V$ be a complex Weil representation with finite inertia image and no monodromy. Its ordinary Artin conductor is
\[
 a(V)=\dim V-\dim V^{I_F}+\Sw(V).
 \tag{22.1}\label{eq:artin}
\]
We use $a$ for the ordinary conductor throughout this part. The Swan conductor, an epsilon exponent for a specified additive character, and the self-pair conductor of a representation of $\GL_n(F)$ are named separately.

Take $\psi_0$ trivial on $\OO_F$ and nontrivial on $\pp_F^{-1}$. Set $\psi_1(x)=\psi_0(\varpi^{-1}x)$, so $\psi_1$ is trivial on $\pp_F$ and nontrivial on $\OO_F$. Under the usual change-of-character convention, if the exponent is written as a power of $q^{1/2-s}$, then
\[
 f_{\psi_0}(V)=a(V),\qquad f_{\psi_1}(V)=a(V)-\dim V.
 \tag{22.2}\label{eq:psi}
\]
The determinant factor in the change of additive character is independent of $s$ and does not alter this exponent shift. Equation \eqref{eq:psi} is a convention translation, not an assertion that every epsilon constant is determined by its exponent.

For a wild-irreducible source $\rho$ of dimension nine with Swan conductor $24$, the audited self-pair calculation below gives the following table.
\begin{center}\begin{tabular}{lrrrrr}\toprule
Object & Dimension & $I_F$-fixed & Swan & Ordinary $a$ & $f_{\psi_1}$\\\midrule
$\rho$ & 9 & 0 & 24 & 33 & 24\\
$\End\rho$ & 81 & 1 & 152 & 232 & 151\\
$\Ad\rho$ & 80 & 0 & 152 & 232 & 152\\\bottomrule
\end{tabular}\end{center}
Here $\Ad\rho$ denotes the trace-zero summand. Over the complex coefficient field, $\End\rho=\C\oplus\Ad\rho$. The scalar summand is unramified, so it contributes no ordinary conductor. It does contribute one to the dimension shift in \eqref{eq:psi}.

Thus the number $152$ has two legitimate appearances, as the Swan conductor of $\End\rho$ and as the $\psi_1$ exponent of its trace-zero summand. It is \emph{not} the $\psi_1$ exponent of all of $\End\rho$, which is $151$. The formal-degree parameter $\delta$ appearing later also equals $152$ in this particular torsion-one comparison. That numerical equality does not make the three definitions interchangeable.

\source{The ordinary conductor identity and additive-character normalization underlie [Q14]. [BH99, \S6.1, Lemma 2] uses the character trivial on $\pp_F$. [AP, \S4 and Theorem 6.9] uses the ordinary pair conductor. The table is an explicit reconciliation of these conventions using the source calculation in Section~\ref{sec:ranknine}.}

\subsection{The principal-order grade calculation}
Let $\A\subset M_n(F)$ be a principal hereditary order of period $e$, with radical $\Pp$. Then $e\mid n$. A principal lattice chain has $e$ successive residue pieces, each of dimension $m=n/e$ over $k_F$. The homogeneous grade of an endomorphism is the direct sum of $e$ matrix blocks of size $m$ by $m$. Consequently
\[
 \dim_{k_F}\Pp^j/\Pp^{j+1}=e m^2=\frac{n^2}{e},
 \qquad [\Pp^{-r}:\A]=q^{r n^2/e} \quad(r>0).
 \tag{22.3}\label{eq:grade}
\]
The period here belongs to the actual principal order carrying the stratum. It cannot be borrowed from a more convenient unrelated lattice chain.

For an actual cuspidal representation containing a positive-level simple stratum $[\A,r,0,\beta]$ to which the hypotheses of [BH99, \S6.1] apply, its Lemma 2 gives
\[
 f_{\psi_1}(\pi)=\frac{nr}{e},\qquad
 a(\pi)=n+\frac{nr}{e}.
 \tag{22.4}\label{eq:principal-conductor}
\]
The first number is the native epsilon exponent in that source. In the present wild cuspidal comparison it equals the standard Swan conductor. A lower critical level in a defining sequence is not the leading level $r$ in this formula.

\begin{proposition}[Necessary period divisibility]\label{prop:period}
Suppose an actual type satisfying \eqref{eq:principal-conductor} is to match a source of rank $n$ and Swan conductor $w>0$. Its principal period and leading level necessarily satisfy
\[
 \frac{n}{\gcd(n,w)}\mid e\mid n,
 \qquad r=\frac{ew}{n}.
 \tag{22.5}\label{eq:period}
\]
These are necessary numerical conditions, not sufficient existence conditions for a simple stratum or a prescribed source.
\end{proposition}
\begin{proof}
The principal condition gives $e\mid n$. Matching \eqref{eq:principal-conductor} gives $nr=ew$, so $n\mid ew$. Writing $n=gn'$ and $w=gw'$ with $\gcd(n',w')=1$ yields $n'\mid e$. Conversely these divisibilities make $ew/n$ an integer. They say nothing about purity, intertwining, or character occurrence.
\end{proof}

For $(n,w)=(9,24)$, the possibilities are $(e,r)=(3,8)$ and $(9,24)$. A maximal-order construction of period one cannot have this conductor. The earlier unramified depth-two candidate has Swan conductor $18$ and ordinary conductor $27$, not $24$ and $33$. Replacing the leading level eight by a critical level four in the period-three formula would create a different error rather than repair the comparison.

\subsection{A self-twist check that survives conductor matching}
A second necessary invariant is the torsion number, the order of the group of unramified self twists. For a maximal simple type, the exact formula in [BK, 6.2.5] identifies that number as $n/e(K/F)$ for the actual parameter field $K$.

\begin{lemma}[Wild irreducibility excludes nontrivial unramified self twists]\label{lem:source-twist}
If $\rho|_{P_F}$ is irreducible and an unramified character $\chi$ satisfies $\rho\otimes\chi\simeq\rho$, then $\chi=1$.
\end{lemma}
\begin{proof}
An intertwiner restricts to an endomorphism of the irreducible $P_F$ representation, because $\chi$ is trivial on $P_F$. It is therefore a nonzero scalar. Applying the intertwining identity to every Weil element forces $\chi$ to be one.
\end{proof}

A period-three candidate may pass the first conductor test and still fail this second test if its actual maximal parameter field has residue degree three. That candidate has torsion number three, while the prescribed source has torsion number one. This is why the conductor equation does not complete the arithmetic attachment.

\section{The rank-nine self-pair filter}\label{sec:ranknine}
\subsection{The source invariant is computed on the unchanged source}
Retain the original source over $F=\F_3((t))$,
\[
 \rho=A_1\otimes A_2,
\]
where $A_h$ is the rank-three Heisenberg representation attached to the projective equation $x_h^9+x_h=t^{-h}$, for $h=1,2$. The independent quotient and wild-irreducibility hypotheses are those of the original source fixture, not newly imposed conditions on a replacement source.

After the recorded unramified splitting extension, take a nonzero constant $v$ satisfying $v^9=-v$. Then
\[
 y_v=v^3x_h^3+vx_h,
 \qquad y_v^3-y_v=-v t^{-h}.
 \tag{23.1}\label{eq:aschar}
\]
The cancellation follows by substituting $v^9=-v$ and $x_h^9+x_h=t^{-h}$. As $v$ varies, these give the eight distinct nontrivial characters of the nine-element translation quotient. A nonzero polynomial of degree three cannot vanish on all nine translations, so the corresponding additive functional is nonzero. Distinct $v$ give distinct functionals by the same argument. Since each pole order $h$ is prime to three, the inherited Artin--Schreier ramification calculation gives break $h$.

\begin{proposition}[Self-pair conductor of the retained source]\label{prop:end-source}
Under the original independence and wild-irreducibility hypotheses,
\[
 \Sw(\End\rho)=152,
 \qquad a(\End\rho)=232.
 \tag{23.2}\label{eq:source232}
\]
\end{proposition}
\begin{proof}
Conjugation on the nine Weyl operators for $A_h$ gives each character of its projective quotient once. Hence $\End A_h$ is the sum of the trivial character and eight characters of break $h$ after the unramified extension. In $\End A_1\otimes\End A_2$, there is one trivial character, eight characters of break one, and $9\cdot8=72$ characters of break two. Independence retains the distinct quotient characters. Products of characters of different breaks have the larger break. The Swan conductor is therefore $8+2\cdot72=152$.

Unramified restriction preserves the Swan conductor. By wild irreducibility, the inertia-fixed endomorphisms are exactly the scalars. Thus $\dim(\End\rho)^{I_F}=1$ over $F$ itself. Equation~\eqref{eq:artin} gives $81-1+152=232$. No assumption that Frobenius fixes each individual quotient character is needed.
\end{proof}

\source{[Q13] supplies the source and the ramification inputs. The displayed Artin--Schreier cancellation, character separation and conductor addition are expanded here. The new exact suite constructs the nine-element constant subspace inside $\F_{81}$ and checks all its additive character values. This finite check does not reprove local ramification theory.}

\subsection{The actual-order certificate cannot be omitted}
The proposed rank-nine family has a complete principal order $\A$ of period nine, radical $\Pp$, and an actual final parameter field $K/F$ that is totally ramified of degree nine. An intermediate field $E$ has degree three. Write $\mathfrak Q$ for the radical of the induced order in the $E$-centralizer. The retained candidate certificate specifies leading level $24$, lower even critical level $c=2a$, and actual simple-character orders
\[
 \mathfrak J=\OO_K+\mathfrak Q^a+\Pp^{12},
 \qquad
 \mathfrak H=\OO_K+\mathfrak Q^{a+1}+\Pp^{13}.
 \tag{23.3}\label{eq:rank9orders}
\]
These equations are substantive construction inputs. The mere pair of numbers $(24,c)$ does not imply them.

The certificate also ensures that the $K$ contribution adds no unexpected graded direction. In the notation of [Q13], minimality of the leading approximation $\gamma$ and the correction $d\in\Pp^{-c}$ give
\[
 [\gamma,\OO_K]\subseteq\Pp^{-c},
 \qquad
 \OO_K\subseteq\mathfrak B+\Pp^{24-c}.
 \tag{23.4}\label{eq:compat}
\]
For $c=4$ and $c=10$, the last exponent exceeds thirteen. Through the grades used in \eqref{eq:rank9orders}, the $K$-grade therefore lies inside the three-dimensional centralizer grade. Its own dimension is one. This is an actual commutator and approximation assertion, not an integral splitting of a wildly ramified bimodule.

Let $I^+=1+\Pp$. The successive full grades have dimension nine, the intermediate grades dimension three, and the $K$-grades dimension one. It follows that
\begin{align}
 I&=\log_q[I^+:J^1]=8(a-1)+6(12-a)=2a+64,\tag{23.4a}\label{eq:rank9I}\\
 D&=\log_q[J^1:H^1]=2+6=8,
 \qquad \dim\eta=q^4.\tag{23.4b}\label{eq:rank9D}
\end{align}
Since $K/F$ is totally ramified of full degree, the finite residue factor is a character and the torsion number is one.

\subsection{A formal-degree calculation with its measure fixed}
Normalize Haar measure by $\operatorname{vol}(\GL_9(\OO_F))=1$ and central measure by $\operatorname{vol}(\OO_F^\times)=1$. Then
\[
 \operatorname{vol}(I^+)=\prod_{j=1}^{9}(q^j-1)^{-1},
 \qquad
 \operatorname{vol}(J)=(q-1)q^{-I}\operatorname{vol}(I^+).
\]
The compact-induction degree formula and the Steinberg degree formula give
\[
 d(\pi)=\frac{q^4}{9\operatorname{vol}(J)},
 \qquad
 d(\mathrm{St}_9)=\frac1{9}\prod_{j=1}^{8}(q^j-1),
\]
so that
\[
 \frac{d(\pi)}{d(\mathrm{St}_9)}
 =\frac{q^9-1}{q-1}\,q^{2a+68}.
 \tag{23.5}\label{eq:degree9}
\]
The factor nine comes from the actual noncompact extension $[K^\times J:F^\times J]$. Omitting it before comparing the measures changes the calculation. An unramified twist may be used to unitarize solely for the formal-degree comparison. That operation does not replace the prescribed noncompact scalar in the construction, and it leaves the self-pair unchanged.

The primary all-characteristic formulas [AP, Theorems 6.5 and 6.9] use the torsion number $r_t$ and an auxiliary parameter $\delta$:
\[
 \frac{d(\pi)}{d(\mathrm{St}_n)}
 =r_t\frac{q^n-1}{q^{r_t}-1}
 q^{(r_t-n+\delta)/2},
 \qquad
 a(\pi^\vee\times\pi)=\delta+n^2-r_t.
 \tag{23.6}\label{eq:AP}
\]
These are the named classical inputs here, not a conjectural formal-degree identity. With $n=9$, $r_t=1$ and \eqref{eq:degree9}, they give
\[
 \delta=4a+144=2c+144,
 \qquad a(\pi^\vee\times\pi)=4a+224=2c+224.
 \tag{23.7}\label{eq:pair9}
\]

\begin{center}\begin{tabular}{lrrrrr}\toprule
Certified profile & $c$ & $I$ & $D$ & $\delta$ & Self-pair $a$\\\midrule
First repair & 10 & 74 & 8 & 164 & 244\\
Deeper repair & 4 & 68 & 8 & 152 & 232\\
Retained source & --- & --- & --- & 152 & 232\\\bottomrule
\end{tabular}\end{center}

\begin{corollary}[A necessary selection inside this certified family]\label{cor:c4}
If a member of the actual-order family \eqref{eq:rank9orders} is to match the retained source under a correspondence preserving the ordinary self-pair conductor, then $c=4$. The critical-ten member is excluded. The critical-four member passes this invariant, but the comparison does not prove its source identification.
\end{corollary}
\begin{proof}
Equate \eqref{eq:source232} and \eqref{eq:pair9}. The equation $2c+224=232$ has the unique solution $c=4$. Equality of this one invariant supplies no character-occurrence or full pair-factor identity.
\end{proof}

\source{[Q13] provides the conditional actual-order certificate and calculation. [AP, formulas (17), (29), Theorems 6.5, 6.9 and Remark 6.7] were checked at their stated measures and field scope. This block does not newly construct the proposed corrected stratum. It preserves the requirement to prove its actual order equations before applying the numerical filter.}

\section{Why pair conductors do not determine quotient capacity}\label{sec:profiles}
\subsection{The parity identity behind the graded calculation}
For the totally ramified complete-order sector, write the successive graded centralizer dimensions as
\[
 n=b_0>b_1>\cdots>b_r=1,
 \qquad \Delta_j=b_{j-1}-b_j.
\]
Assume the actual defining sequence and its order compatibility give descending positive levels $\nu_j$ and the grade-count formulas
\[
 I=\sum_{j=1}^r\Delta_j\left\lceil\frac{\nu_j}{2}\right\rceil-(n-1),
 \qquad
 D=\sum_{j=1}^r\Delta_j\,\mathbf1_{2\mid\nu_j}.
 \tag{24.1}\label{eq:profile}
\]
Here $I=\log_q[I^+:J^1]$, while $D=\dim_{k_F}(J^1/H^1)$. The different floor and ceiling thresholds are essential. In this torsion-one, full-degree setting, the same formal-degree calculation gives
\[
 a(\pi^\vee\times\pi)=n^2-1+(n-1)+2I+D.
 \tag{24.2}\label{eq:pair-profile}
\]
As in the rank-nine calculation, the formula is conditional on actual orders. A list of integers is not a certificate of their arithmetic existence.

\begin{proposition}[Parity cancellation]\label{prop:parity}
Under \eqref{eq:profile},
\[
 (n-1)+2I+D=\sum_{j=1}^r\Delta_j\nu_j.
\]
Consequently
\[
 a(\pi^\vee\times\pi)=n^2-1+\sum_{j=1}^r\Delta_j\nu_j.
 \tag{24.3}\label{eq:pair-weight}
\]
The weighted conductor expression does not retain $D$ as a separate invariant.
\end{proposition}
\begin{proof}
For each integer $\nu$,
\[
 2\left\lceil\frac{\nu}{2}\right\rceil+
 \mathbf1_{2\mid\nu}=\nu+1.
\]
Insert this identity into $2I+D$. Since $\sum_j\Delta_j=n-1$, the constant terms reduce to $-(n-1)$. Adding $n-1$ proves the result.
\end{proof}

This cancellation explains a concrete limitation of the earlier numerical tests. Changing the parity distribution can change the Heisenberg quotient while preserving the weighted sum in \eqref{eq:pair-weight}. More conductor precision does not recover information that the formula has already summed away.

\subsection{The full dyadic profile}
In the inherited full dyadic construction, let $n=2^d$ and let $h_1<\cdots<h_d$ be positive odd integers satisfying its separation requirements. The consecutive profile has
\[
 b_j=2^{d-j},\qquad
 \nu_j=2^{d-j}\,3h_{d-j+1}
 \quad(1\le j\le d).
 \tag{24.4}\label{eq:dyadic-profile}
\]
All levels except the final one are even. Therefore
\[
 D=2^d-2.
 \tag{24.5}\label{eq:dyadic-D}
\]
A source consisting of $d$ nondegenerate binary symplectic planes has dimension $2d$. The original maximal-type quotient is too small for $d=1,2$, has exactly the needed dimension for $d=3$, and has strict surplus for $d\ge4$.

\begin{center}\begin{tabular}{rrrrl}\toprule
$d$ & Arithmetic rank & Source dimension & Host $D$ & Original quotient\\\midrule
1 & 2 & 2 & 0 & Insufficient\\
2 & 4 & 4 & 2 & Insufficient\\
3 & 8 & 6 & 6 & Equal capacity\\
4 & 16 & 8 & 14 & Strict surplus\\
5 & 32 & 10 & 30 & Strict surplus\\\bottomrule
\end{tabular}\end{center}

The rank-two and rank-four enlarged arithmetic carriers in Part I are not counterexamples to the first two rows. Those repairs change the assembly carrier while recovering the same inducing representation. The table concerns the original specified maximal-type quotient. At rank eight, dimension ceases to be the obstruction, but an independently fixed quadratic marking can still obstruct the attachment.

\subsection{Two rank-eight profiles with the same conductor data}
The retained dyadic capacity note contains a sharper example. In the full rank-eight profile, use
\[
 (b_0,b_1,b_2,b_3)=(8,4,2,1),
 \qquad (\nu_1,\nu_2,\nu_3)=(12h_3,6h_2,3h_1).
\]
A second constructed profile skips the intermediate quartic field and has
\[
 (b_0,b_1,b_2)=(8,4,1),
 \qquad (\nu_1,\nu_2)=(12h_3,c),\qquad c=h_1+4h_2.
 \tag{24.6}\label{eq:skip}
\]
The integer $c$ is odd and $2c<12h_3$. Its last jump therefore contributes nothing to $D$. The full profile has $D=6$, while the skipped profile has $D=4$.

Both weighted sums in \eqref{eq:pair-weight} equal
\[
 48h_3+12h_2+3h_1.
 \tag{24.7}\label{eq:sameweight}
\]
Indeed, the full profile has jump sizes $4,2,1$, and the skipped profile has jump sizes $4,3$. The identity $3(h_1+4h_2)=3h_1+12h_2$ supplies the equality.

For $(h_1,h_2,h_3)=(1,3,5)$, the distinction becomes completely explicit.
\begin{center}\begin{tabular}{lrr}\toprule
Quantity & Full profile & Skipped profile\\\midrule
Levels & $(60,18,3)$ & $(60,13)$\\
Principal period & 8 & 8\\
Standard Swan conductor & 60 & 60\\
Standard ordinary conductor & 68 & 68\\
Torsion number & 1 & 1\\
$I$ & 133 & 134\\
$D$ & 6 & 4\\
Heisenberg dimension over $q=2$ & 8 & 4\\
$\delta$ & 279 & 279\\
Ordinary self-pair conductor & 342 & 342\\\bottomrule
\end{tabular}\end{center}
The compensating change is $2\cdot133+6=2\cdot134+4$. The three necessary invariants---standard conductor, torsion number, and self-pair conductor---cannot force the full six-dimensional quotient.

\subsection{The skipped profile is an arithmetic claim, not an integer exercise}
The inherited construction of \eqref{eq:skip} works over $\F_2((t))$. Put $q_0=3h_3$ and take $\alpha^2+\alpha=t^{-q_0}$. A uniformizer $z$ of the quadratic field satisfies
\[
 z=t^{(q_0+1)/2}\alpha,
 \qquad z^2+t^{(q_0+1)/2}z+t=0.
\]
Over that field, the polynomial $X^4+zX+z$ is separable and Eisenstein. If $e$ is a root, let $C$ be multiplication by $e^{-c}$ in the quartic centralizer. The source note chooses a filtered corestriction lift using $s(P_0)=t/z^2$ and takes a correction proportional to $(z^2/t)CP_0$. The actual final parameter is the literal corrected matrix, not an assumed compositum of the approximation fields.

The conclusion that this correction supplies the claimed defining sequence uses the retained refinement argument and its strict inequalities. This publication block has audited the conductor and capacity consequence of that argument. It has not freshly rerun a skipped-profile matrix constructor or recertified every step of its arithmetic field proof. Those statuses are not replaced by the exact integer tests.

Nor is this example a counterexample to H5. Full local pair factors are substantially more information than the three numbers above. The example shows that those three numbers alone do not identify the needed native quotient, and that a source-defined candidate cannot be declared to contain a prescribed type merely because it passes them.

\source{[D167] supplies the actual profile formulas and the skipped-degree construction. Proposition~\ref{prop:parity} and the numerical comparison are fully derived here. The exact suite checks the example and 120 positive odd triples, as well as general parity identities. The underlying arithmetic construction retains its inherited ordinary-proof status.}

\section{The source-defined arithmetic candidate and its inputs}\label{sec:induction}
\subsection{Do not identify two candidates by notation}
The direct matrix recursion and the classical arithmetic-induction route start with different data. The first produces an actual type satisfying a restricted construction contract. The second starts from a complete source character and uses classical induction and descent to produce a candidate. We denote the latter by $\pi_{\mathrm{BH}}(\rho)$ in this section. This notation does not identify it with any independently constructed compact type.

The audited source-defined branch is over $\F_3((t))$ and uses the original rank-nine source. It is distinct from the characteristic-zero continuation113 candidate discussed in Part II. Its existence through classical inputs does not prove the corrected-character occurrence question for that other candidate, nor does it place a prescribed native packet in the current one.

\subsection{The degree-nine field does not initially have a cubic tower}
Write $A=t^{-1}$, and retain the projective coordinates
\[
 x^9+x=A,\qquad z^9+z=A^2.
\]
The source record takes $y=z-x-x^3$. In characteristic three,
\[
 y^9+y=B,\qquad B=A^2-A-A^3.
 \tag{25.1}\label{eq:B}
\]
The inducing field is $M=F(y)$, of degree nine. Its retained Galois-closure analysis has translation subgroup $(C_3)^2$ with a $C_4$ action irreducible over $\F_3$. There is no cubic intermediate field over $F$. Therefore one cannot simply write two cubic automorphic-induction arrows over $F$ without changing the base and proving the resulting descent.

Pass to the unramified extension $K=\F_{81}((t))$. The constants now split the translation action. Choose $a\ne0$ with $a^9=-a$ and $i$ with $i^2=-1$, and define
\[
 r=\frac{y^3-a^2y}{ia^3},\qquad s=\frac ya.
\]
Then
\[
 r^3-r=\frac{B}{ia},\qquad s^3-s=ir.
 \tag{25.2}\label{eq:tower}
\]
To check the first equality, use $a^8=-1$ to cancel the $y^3$ terms after cubing. The second follows immediately from the definition of $r$. These identities exhibit the cyclic cubic tower over $K$ used by the classical route. The degree and regularity assertions remain part of the complete source fixture.

The inducing character is also complete data. The inherited record includes its ramified Artin--Schreier class, its unit values, and an actual unramified factor. Replacing it by the displayed pole order alone loses information needed for the resulting central character and noncompact extension. The noted different exponent $24$ and character conductor $9$ give the expected induced ordinary conductor $33$, but do not determine that full character.

\subsection{Induction, invariance, and descent are separate gates}
The construction first applies cyclic automorphic induction at the two cubic stages over $K$. Cuspidality requires the recorded regularity at each relevant step. It then invokes the classical independence of the resulting monomial construction from an admissible cyclic tower. This independence is a substantial theorem. It is what allows a Galois conjugate of the chosen tower to be compared with the original construction. The fact that the source descends to $F$ does not, by itself, make each arbitrarily chosen tower Frobenius invariant.

Only after this invariance has been established does the tame descent theorem apply. Its fibre is described by the unramified characters associated to $K/F$, here of order dividing four. The specified determinant of $\rho$ selects the desired central character.

\begin{lemma}[Unique fixed-center member in a coprime descent fibre]\label{lem:descent}
Let a descent fibre be a torsor for a cyclic character group $X$ of order $m$. Suppose twisting a rank-$n$ member by $\alpha\in X$ changes its central character by $\alpha^n$. If $\gcd(n,m)=1$, then any attainable target central character selects exactly one member of the fibre.
\end{lemma}
\begin{proof}
The map $\alpha\mapsto\alpha^n$ is a bijection of $X$ because multiplication by $n$ is invertible modulo $m$. Given one member and an attainable ratio of central characters, there is a unique twist realizing that ratio.
\end{proof}

For $n=9$ and $m=4$, the required power map is invertible. This explains the central normalization rather than silently selecting a preferred descent. It differs from the fixed-center cubic ambiguity in continuation113, where the relevant power map is not coprime. The new finite tests include both the unique and the nonunique cases.

\subsection{What the primary-source audit does and does not establish}
The selected audit uses three layers of classical material.

\textbf{The 1999 monomial and tame-lifting statements.} In [BH99], the definitions and properties of tame lifting in Section 1.3 identify the descent fibre, twisting and central-character behavior. Section 2.3 gives the characteristic-zero monomial construction, including standard epsilon factors and pair $L$-functions and conductors. Its proof uses independence of the cyclic construction. Section 1.7 supplies the tame transfer rule for the relevant pair-conductor exponent. These statements were read at their exact pages.

\textbf{The positive-characteristic extension.} The original Section 2.3 is explicitly characteristic zero. The characteristic note on page eight of [BH13] states that the relevant earlier results extend to positive characteristic through the base-change and automorphic-induction work of [HL11]. That note was checked directly. The present audit therefore retains the positive-characteristic induction theorem as a named external input, rather than erasing the original hypothesis.

\textbf{The limit of fresh source access.} The complete [HL11] PDF could not be retrieved in this block. Its bibliographic record was accessible, but repeated full-text requests failed. The detailed statements III.2 and III.3, including the restricted pair statement used in [Q15], retain their inherited source-audit status. They are not labeled freshly reread here. This is a narrower limitation than saying the entire arithmetic route was unaudited.

There are two particularly important exclusions. First, the effective construction in [BH13] treats the totally wildly ramified correspondence as given at its starting point. That paper's effective theorem is not used here as an independent solution of the same totally wild input. Its characteristic note is the relevant source interface. Second, the equality of pair conductors in [BH99, 2.3(vii)] is not equality of full pair epsilon factors. The introduction expressly distinguishes these strengths.

\subsection{The resulting candidate still has a native attachment obligation}
Under the recorded source and classical hypotheses, [Q15] obtains a cuspidal candidate with prescribed central character, torsion number one, standard Swan conductor $24$, ordinary conductor $33$, and ordinary self-pair conductor $232$. The characterizing information required by the paper's H5 contract is not replaced by that list.

In particular, one cannot compare an essentially discrete source with a reducible principal series outside the scope of a pair theorem and then divide a pair conductor by nine to manufacture the standard conductor. Nor can a first-order calculation of a parameter on deep field units establish occurrence of a full simple character. The original unit domain and noncompact extension remain part of that question.

The next useful arithmetic step is therefore not another calculation of the same three integers. It is to construct or identify the actual native type and its complete character in this fixed candidate, and to compare its marked data with the supplied source. That is the boundary the finite results in the following sections are designed to expose.

\source{[Q15] is the inherited source-defined constructor audit. [BH99, printed pp. 15, 18, 23--25, 53--54] and [BH13, p. 8] were freshly checked. [HL11] remains a named classical input with an explicit fresh-access limitation. No source-defined candidate is substituted for a separately fixed type candidate without a proof of their identification.}

\section{A full dyadic extension remembers its squares}\label{sec:extensions}
\subsection{The extension criterion with a fixed center}
Let $V$ be a finite-dimensional vector space over $\F_2$, and consider central extensions
\[
 1\longrightarrow Z\longrightarrow E_s\longrightarrow V\longrightarrow1,
 \qquad
 1\longrightarrow Z\longrightarrow E_t\longrightarrow W\longrightarrow1,
\]
where $Z$ is the same abelian group in both extensions. In applying this to a represented scalar pushout, the represented center and any already factored character kernel must be recorded. A pushout cannot recover invisible elements that have been discarded.

Fix a linear isomorphism $f:V\to W$. Choose an ordered basis $(e_i)$ of $V$, lifts $a_i$ in $E_s$, and lifts $b_i$ of $f(e_i)$ in $E_t$. Put $s_i=a_i^2$ and $t_i=b_i^2$ in $Z$. Write $Z^2=\{z^2:z\in Z\}$.

\begin{proposition}[Center-fixing extension criterion]\label{prop:ext}
There exists an isomorphism $F:E_s\to E_t$ fixing $Z$ pointwise and lifting $f$ if and only if
\[
 [a_i,a_j]=[b_i,b_j]\quad\hbox{for every }i,j,
 \qquad s_i t_i^{-1}\in Z^2\quad\hbox{for every }i.
 \tag{26.1}\label{eq:extcrit}
\]
If the generating lifts themselves must satisfy $F(a_i)=b_i$, then their squares must agree exactly, not merely modulo $Z^2$.
\end{proposition}
\begin{proof}
Any such isomorphism has $F(a_i)=\lambda_i b_i$ for some $\lambda_i\in Z$. Preservation of squares gives $s_i=\lambda_i^2 t_i$, while preservation of commutators gives the first condition. Conversely choose $\lambda_i$ with these squares. Every group element has a unique normal form
\[
 z a_1^{v_1}\cdots a_r^{v_r},\qquad z\in Z,\quad v_i\in\{0,1\}.
\]
Define its image by replacing $a_i$ with $\lambda_i b_i$. Normal-form multiplication uses only central elements, the displayed squares, and the pairwise commutators. The conditions ensure that each multiplication relation is preserved. The inverse construction proves bijectivity. Taking every $\lambda_i=1$ gives the final assertion.
\end{proof}

For $Z=\mu_2=\{1,-1\}$, each square is independent of the chosen lift and defines a quadratic refinement
\[
 a_v^2=(-1)^{q(v)},\qquad
 q(v+w)=q(v)+q(w)+B(v,w),
 \tag{26.2}\label{eq:quadratic}
\]
where $(-1)^{B(v,w)}$ is the commutator. A center-fixing extension isomorphism over a specified symplectic map $f$ therefore requires $q_t\circ f=q_s$.

For $Z=\mu_4$, multiplying a lift by $i$ changes its square by $-1$. The unmarked square classes are consequently weaker. For $Z=\C^\times$, every square ratio has a square root and the unmarked criterion reduces to matching commutators. Neither statement deletes a separately marked finite subgroup or a prescribed family of lifts.

\subsection{Dihedral and quaternion data inside the same scalar pushout}
Consider the matrices
\[
 X=\begin{pmatrix}0&1\\1&0\end{pmatrix},\qquad
 Z_0=\begin{pmatrix}1&0\\0&-1\end{pmatrix}.
\]
They satisfy $X^2=Z_0^2=1$ and $XZ_0=-Z_0X$. The group $H_D=\langle X,Z_0\rangle$ is a dihedral group of order eight. The group $H_Q=\langle iX,iZ_0\rangle$ is a quaternion group of order eight. On their common two-dimensional quotient, their refinements are
\[
 q_D(x,y)=xy,\qquad q_Q(x,y)=xy+x+y.
 \tag{26.3}\label{eq:DQ}
\]

\begin{center}\begin{tabular}{lrrrr}\toprule
Marked extension & $q(e)$ & $q(f)$ & $q(e+f)$ & $\sum_v(-1)^{q(v)}$\\\midrule
$H_D$ & 0 & 0 & 1 & $+2$\\
$H_Q$ & 1 & 1 & 1 & $-2$\\\bottomrule
\end{tabular}\end{center}

Both scalar pushouts identify with
\[
 P=\{zX^aZ_0^b:z\in\C^\times,\ a,b\in\{0,1\}\}.
\]
The marked subgroups remain different. In $H_D$, five nonidentity elements have order two and two have order four. In $H_Q$, only the central nonidentity element has order two and six have order four. In particular, no center-fixing automorphism of $P$ carries one marked subgroup onto the other.

On the other hand, a quaternion source does embed into $P$ through $iX,iZ_0$. Its squares and negative Gauss sign are retained. What fails is identifying it with the \emph{independently chosen} dihedral subgroup $H_D$. This example separates a successful transported marking from a false assertion about a fixed native marking.

\subsection{A full marked section contains more than a square table}
Suppose a normalized section $\sigma:V\to E$ has product cocycle
\[
 \sigma(v)\sigma(w)=c(v,w)\sigma(v+w).
\]
If an isomorphism must send each specified section element to a specified target section element, the two cocycles must agree after transport on \emph{every} pair. Equality of polar forms and square tables only identifies their extension class, not all those pointwise choices.

For example, start with the ordered section $\sigma(x,y)=X^xZ_0^y$ of $H_D$. Change only the value at $e+f$ by a central sign. Every square remains unchanged because $(-1)^2=1$, but the product of the specified lifts at $e$ and $f$ now differs by a sign from the specified lift at $e+f$. An isomorphism fixing the two generating lifts cannot also implement this changed third assignment. The new tests retain this counterexample rather than treating it as a harmless convention change.

\subsection{The valid scalar-pushout attachment contract}
There is nevertheless a positive attachment statement in the restricted sector already used by Paper XI. Suppose an actual type, its full character and its Heisenberg representation $\eta_W$ are supplied. Suppose the graph has one vertex, the prescribed stabilizer is trivial, and the source scalar pushout has the same faithfully represented center and a nondegenerate quotient $V$ with $\dim V\le\dim W$. No conflicting host-native finite marking is independently fixed.

Choose a symplectic embedding $f:V\to W$. Proposition~\ref{prop:ext} gives a scalar-center-fixing isomorphism of the source pushout with the restriction over $f(V)$. Carry every supplied finite subgroup, every lift, and every section value through this \emph{same} isomorphism. Since it is a homomorphism fixing scalars, it preserves all squares, commutators and section products, not just their classes.

Writing $W=f(V)\perp U$, the finite Heisenberg representation factors as $\eta_V\otimes\eta_U$. A projector onto a line in the second factor gives the source Heisenberg sector inside the actual image algebra. The inherited recovery theorem then returns the unchanged full inducing representation and candidate, under its stated recovery hypotheses. There is no stabilizer equation because the stipulated stabilizer is trivial.

This is an attachment into the scalar extension of the actual type. It is not a claim that every source lift is an unscaled arithmetic matrix, nor that the transported subgroup equals an independently specified canonical host subgroup. A package that records only the unmarked ambient scalar extension loses precisely the information displayed by \eqref{eq:DQ}.

\source{[Q16] contains the complete marked scalar-pushout contract. The normal-form proof, the dihedral/quaternion calculations and the section counterexample are explicit finite arguments. Their application uses the already supplied actual type and inherited recovery theorem. This section does not construct those arithmetic inputs.}

\section{The sharp boundary for a fixed native refinement}\label{sec:capacity}
\subsection{Quadratic type and its Gauss sign}
Let $(V,B)$ be a nondegenerate alternating space over $\F_2$ of dimension $2s$, and let $q$ be a quadratic refinement. In a symplectic basis $(e_j,f_j)$,
\[
 q(x_1,y_1,\ldots,x_s,y_s)
 =\sum_{j=1}^s\bigl(x_jy_j+q(e_j)x_j+q(f_j)y_j\bigr).
\]
Its Gauss sum factors into plane sums:
\[
 G(q)=\sum_{v\in V}(-1)^{q(v)}
 =2^s(-1)^{\sum_j q(e_j)q(f_j)}.
 \tag{27.1}\label{eq:gauss}
\]
Define $\Arf(q)\in\F_2$ by the sign in \eqref{eq:gauss}. The sum is intrinsic, so this definition is independent of the symplectic basis.

\begin{lemma}[Classification in the required finite setting]\label{lem:arfclass}
Two nondegenerate quadratic spaces over $\F_2$ are isometric if and only if they have the same dimension and Arf invariant. Both Arf types occur in every positive even dimension.
\end{lemma}
\begin{proof}
The coordinate expression decomposes a space into nondegenerate planes. On a plane, the three refinements with Arf invariant zero are isometric to $H(x,y)=xy$, while the unique remaining type is $A(x,y)=xy+x+y$. This can be checked on its three nonzero vectors.

Two anisotropic planes can be replaced by two hyperbolic planes. In $A\perp A$, the vector $u=e_1+e_2$ is nonzero and isotropic. Choose $v$ with $B(u,v)=1$. Replacing $v$ by $v+q(v)u$ makes it isotropic too, so $u,v$ span a hyperbolic plane. Its nondegenerate orthogonal complement is a plane whose Gauss sign is positive by multiplicativity of \eqref{eq:gauss}. It too is hyperbolic. Pairing all anisotropic planes leaves either only hyperbolic planes or one anisotropic plane. The parity of the remaining anisotropic planes is exactly the Arf invariant. These normal forms prove classification and existence.
\end{proof}

In dimension six, \eqref{eq:gauss} gives $G(q)=+8$ or $-8$. The numbers of zero values are respectively
\[
 \#q^{-1}(0)=\frac{64+8}{2}=36,
 \qquad \#q^{-1}(0)=\frac{64-8}{2}=28.
 \tag{27.2}\label{eq:rank8zeros}
\]
These counts include the zero vector. The new suite verifies both refinements on all 64 vectors. They supply a concrete signature of the distinction that dimension alone misses.

\subsection{One spare plane is enough only when the embedding is free}
\begin{proposition}[Fixed-refinement capacity criterion]\label{prop:capacity}
Let $(V,q_V)$ and $(W,q_W)$ be nondegenerate quadratic spaces of dimensions $2s$ and $2R$. When the embedding may be freely chosen, a quadratic isometric embedding $V\to W$ exists exactly when
\[
 R\ge s,
 \qquad R=s\ \Longrightarrow\ \Arf(q_V)=\Arf(q_W).
 \tag{27.3}\label{eq:capcrit}
\]
If $R>s$, either source Arf type embeds.
\end{proposition}
\begin{proof}
An embedding requires the dimension inequality. At equal dimension it is an isometry, so the Arf invariants must agree. If $R>s$, choose a nondegenerate space $U$ of dimension $2(R-s)$ with
\[
 \Arf(q_U)=\Arf(q_W)+\Arf(q_V).
\]
Both choices are available by Lemma~\ref{lem:arfclass}. The space $V\perp U$ has the same dimension and Arf invariant as $W$, hence is isometric to $W$. Restrict that isometry to $V$. The equal-dimension sufficiency follows from the same classification.
\end{proof}

After the quadratic embedding is chosen, the preimage of its image in an actual $\mu_2$-central extension of $W$ has the source refinement. Proposition~\ref{prop:ext} supplies the extension isomorphism. Exact generating lifts can be matched when their squares and commutators match. Additional prescribed section values retain the stronger product condition from Section~\ref{sec:extensions}.

The spare-plane argument does not permit arbitrary changes to fixed depth placements, transport maps, or stabilizer actions. It proves existence of an otherwise free quadratic embedding. These qualifications are important even when the host has a large dimensional surplus.

\subsection{The exact rank-eight consequence}
At $d=3$ in \eqref{eq:dyadic-D}, the full dyadic quotient and the three-plane source both have dimension six. Suppose the source planes have Arf invariants $\epsilon_1,\epsilon_2,\epsilon_3$, and a particular actual host refinement has Arf invariant $\epsilon_W$. Then the necessary and sufficient test for a \emph{free} quadratic embedding is
\[
 \epsilon_W=\epsilon_1+\epsilon_2+\epsilon_3\quad\hbox{in }\F_2.
 \tag{27.4}\label{eq:rank8arf}
\]
There is no orthogonal complement in which to absorb a mismatch. For $d\ge4$, the full profile has a positive-dimensional complement, so the finite Arf obstruction disappears when the embedding is free and the actual host refinement is supplied.

This does not compute $\epsilon_W$ for a particular arithmetic type. A Heisenberg commutator quotient alone supplies an alternating form. It does not canonically select a finite two-element-center extension or its square refinement. Determining an independently required native refinement and its lift phases for the chosen candidate remains an arithmetic obligation.

Conversely, under the transported scalar-pushout contract of Section~\ref{sec:extensions}, the host refinement is part of the transported output rather than independently prescribed. Every source marking can then be retained in the same scalar extension once the symplectic capacity and recovery hypotheses hold. There is no contradiction between this positive statement and \eqref{eq:rank8arf}. They quantify over different data.

\source{The capacity criterion was recorded in [Q16]. Its finite classification input and proof are made explicit here. The rank-eight deduction combines that criterion with [D167]. The block does not assign a canonical native Arf invariant to an actual arithmetic candidate that has not supplied one.}

\section{An explicit correction for a quadratic marking defect}\label{sec:correction}
\subsection{The discrepancy is a vector, not an unspecified search}
Suppose a symplectic embedding $f:(V,B_V)\to(W,B_W)$ has already been supplied. Put
\[
 q'(v)=q_W(f(v)),\qquad \ell(v)=q'(v)+q_V(v).
\]
The two quadratic refinements have the same polar form, so $\ell$ is linear. Nondegeneracy gives a unique vector $a\in V$ such that
\[
 \ell(v)=B_V(a,v)\quad\hbox{for every }v.
 \tag{28.1}\label{eq:defect}
\]
Finding $a$ is therefore one linear solve over $\F_2$. It is not a search over all symplectic automorphisms.

\begin{lemma}[Arf change under a linear defect]\label{lem:arfshift}
With the above notation,
\[
 \Arf(q')+\Arf(q_V)=q_V(a).
 \tag{28.2}\label{eq:arfshift}
\]
\end{lemma}
\begin{proof}
The refinement identity gives
\[
 q_V(v)+B_V(a,v)=q_V(v+a)+q_V(a).
\]
Translation by $a$ permutes $V$. Summing the signs yields $G(q')=(-1)^{q_V(a)}G(q_V)$. Equation~\eqref{eq:gauss} identifies the two signs.
\end{proof}

Define the transvection
\[
 T_a(v)=v+B_V(a,v)a.
 \tag{28.3}\label{eq:transvection}
\]
It is symplectic and satisfies $T_a^2=1$. For $a=0$ it is the identity.

\begin{proposition}[One-transvection quadratic correction]\label{prop:transvection}
If $q_V(a)=0$, then
\[
 q'\bigl(T_a(v)\bigr)=q_V(v)\quad\hbox{for every }v.
 \tag{28.4}\label{eq:corrected}
\]
Thus $f\circ T_a$ is a quadratic isometric embedding. The condition is equivalent to equality of the Arf invariants of $q'$ and $q_V$.
\end{proposition}
\begin{proof}
Write $b=B_V(a,v)$. Since $B_V(a,a)=0$, one has $B_V(a,T_av)=b$. Also
\[
 q_V(T_av)=q_V(v)+bq_V(a)+b.
\]
Adding the defect term $B_V(a,T_av)=b$ gives
\[
 q'(T_av)=q_V(v)+bq_V(a).
\]
This equals $q_V(v)$ when $q_V(a)=0$. For completeness, expanding the two cross terms in $B_V(T_av,T_aw)$ makes them cancel in characteristic two, so $T_a$ is symplectic. Its square is the identity because the same value $b$ is added twice. Lemma~\ref{lem:arfshift} proves the equivalence of the condition.
\end{proof}

This yields a direct finite correction procedure: compute the defect on a basis, solve for $a$, test $q_V(a)$, and, when it vanishes, compose with $T_a$. The test suite verifies the formula for all linear defects in dimensions two, four and six, against every vector. It also checks the symplectic identities independently.

\subsection{Equal Arf does not validate an unchanged map}
Take $V=\F_2^2$ with $q_V(x,y)=xy$ and $q'(x,y)=xy+y$. Both refinements have Arf invariant zero. The identity map does not preserve the refinement: their values differ at $(0,1)$. The defect vector is $a=(1,0)$, which is isotropic. The transvection $(x,y)\mapsto(x+y,y)$ repairs the map.

The existence of that repair is not a proof that the original identity map worked. When the original embedding is prescribed by an arithmetic placement or a transport diagram, replacing it must itself be permitted and checked. A corrected finite embedding must never be silently substituted into a diagram that required the old map.

\subsection{Preserving a previously correct subspace}
There is a useful relative form. Let $U\subseteq V$ be a subspace on which the defect vanishes. Then $a\in U^\perp$, so $T_a(u)=u$ for every $u\in U$. Under the isotropy condition $q_V(a)=0$, the correction therefore fixes all previously correct vectors in $U$ pointwise.

This statement does not automatically preserve selected lifts above those vectors. It fixes the quotient embedding there. The extension criterion and section products still determine whether the chosen lifted correction respects all prescribed elements. The distinction is small in notation but essential in a full marked packet.

\subsection{A conditional equivariant form}
Suppose a finite group $C$ acts symplectically on $V$ and $W$, preserves $q_V$ and $q_W$, and the supplied $f$ is $C$-equivariant. Then
\[
 \ell(gv)=q_W(f(gv))+q_V(gv)=\ell(v).
\]
Nondegeneracy in \eqref{eq:defect} implies $ga=a$ for every $g\in C$. Consequently $T_a$ commutes with the action. When $q_V(a)=0$, the correction $f\circ T_a$ remains equivariant.

\begin{corollary}[Equivariant finite correction under supplied actions]\label{cor:equiv}
A supplied equivariant symplectic embedding with matching pullback Arf invariant admits the explicit equivariant quadratic correction of Proposition~\ref{prop:transvection}. If its defect already vanishes on a prescribed subspace, the correction fixes that subspace pointwise.
\end{corollary}
\begin{proof}
The invariance computation makes $a$ fixed by $C$. Inserting $ga=a$ into \eqref{eq:transvection} gives $gT_a=T_ag$. The subspace assertion follows from $a\in U^\perp$.
\end{proof}

These are finite statements conditional on actual actions and an initial equivariant embedding. They do not construct such actions inside an arithmetic normalizer. Nor do they show that an arbitrary quadratic embedding provided by a dimension argument can be chosen equivariantly. The relative correction is a way to finish a specified finite step, not a replacement for its input.

\source{This section develops an explicit elementary repair from the linear-defect observation in [Q16]. The Gauss-sum proof and transvection computation are given in full. They are ordinary finite deductions, not a new Lean result or a claim of a new general arithmetic existence theorem.}

\section{Full extension actions and the remaining phase boundary}\label{sec:actions}
\subsection{The discrepancy after a quadratic embedding has been chosen}
Even an equivariant quadratic embedding need not give the required equivariant isomorphism of the full finite extensions with every prescribed marking. This section isolates a conditional repair of the group-action part.

For clarity, take nondegenerate $\mu_2$-central extensions of the same quotient dimension, or restrict a larger target extension to a supplied invariant quadratic subspace. Suppose $C$ acts genuinely on the source and target extensions, fixes the center, and induces the same quotient action after identification. Let $F:E_s\to E_t$ be a center-fixing extension isomorphism inducing that quotient identification. Such an $F$ is supplied by the quadratic extension criterion, but is not initially assumed equivariant.

An automorphism of $E_t$ that fixes its center and quotient is uniquely of the form
\[
 u_\chi(x)=(-1)^{\chi(\bar x)}x,
 \qquad \chi\in\Hom_{\F_2}(W,\F_2).
 \tag{29.1}\label{eq:vertical}
\]
Indeed, the ratio between the image of $x$ and $x$ is a central sign depending only on $\bar x$, and preservation of multiplication makes it additive. Conversely every such character gives an automorphism. These vertical automorphisms form the additive $\F_2$-space $A=\Hom(W,\F_2)$.

Write $\alpha_g$ for the target action and $\beta_g=F\alpha^s_gF^{-1}$ for the transported source action. Their discrepancy is
\[
 \delta_g=\beta_g\alpha_g^{-1}\in A.
\]
The action of $C$ on $A$ is conjugation by $\alpha_g$. Using additive notation in $A$, the genuine group-action equations give
\[
 \delta_{gh}=\delta_g+g\delta_h.
 \tag{29.2}\label{eq:cocycle}
\]
Thus the discrepancy is an actual 1-cocycle, rather than an arbitrary list of mismatched operators.

\subsection{An explicit odd-order averaging repair}
\begin{proposition}[Conditional odd-action correction]\label{prop:odd-action}
Under the preceding hypotheses, suppose $|C|$ is odd. Put
\[
 u=\sum_{h\in C}\delta_h\in A.
 \tag{29.3}\label{eq:average}
\]
Replacing $F$ by $u\circ F$, with $u$ interpreted through \eqref{eq:vertical}, makes the extension isomorphism $C$-equivariant.
\end{proposition}
\begin{proof}
Sum \eqref{eq:cocycle} over $h$. Left multiplication permutes $C$, so
\[
 u=|C|\delta_g+gu=\delta_g+gu
\]
in the characteristic-two space $A$. Thus $\delta_g=u+gu$. Replacing $F$ by $uF$ changes its discrepancy to $u+\delta_g+gu$, which is zero for every $g$.
\end{proof}

The correction fixes the scalar center and the quotient. It may change a chosen lift by a central sign. If all generating lifts or an entire section are independently fixed, that modification may be forbidden. The valid conclusion is existence of an equivariant center-fixing extension isomorphism when these vertical corrections are allowed. It is not an equivariant pointwise identification of every independently prescribed lift.

The exact suite enumerates all cocycles for a $C_3$ action on a two-dimensional $\F_2$ module and checks the averaging formula. It also retains the opposite boundary. For $C_2$ acting trivially on a nonzero $\F_2$ module, a nonzero value of $\delta$ at its generator is a 1-cocycle. Every coboundary $u+gu$ is zero. No correction of this form can remove it. Odd order is therefore a substantive hypothesis, not a cosmetic simplification.

\subsection{What has not been repaired by this finite sum}
The result starts after the arithmetic host, actual finite actions, invariant quotient subspace and an extension isomorphism have been supplied. It does not construct them from the original source diagram. A native-normalizer obstruction can prevent the initial action from existing even when the later finite cocycle would be removable.

There is a further representation-level boundary. An action by automorphisms of a Heisenberg group preserves the isomorphism class of its representation with fixed nontrivial central character. Intertwining operators can therefore be chosen, but their products can carry scalar discrepancies. Passing from those projective operators to the specified honest phase operators is a separate problem. Proposition~\ref{prop:odd-action} removes a discrepancy in the action on the \emph{finite extension group}. It does not prove vanishing of every scalar projective obstruction in a representation, nor does it choose the prescribed Gauss or transport phases.

The original full character domain, its central quasicharacter, the noncompact root, and all specified section products remain part of the arithmetic contract. A finite central sign correction is acceptable only when its interaction with that contract is explicitly checked.

\subsection{A useful order of operations}
The two finite corrections should be used in the following order. First supply an actual invariant host and a quotient-equivariant symplectic embedding. Compare its pulled-back quadratic refinement with the source. When the defect vector is isotropic, use the equivariant transvection to correct the quadratic embedding. Next lift to the finite extensions and compute the resulting action cocycle. When the stipulated hypotheses permit it, use the odd-order sum to correct that action. Finally check every prescribed lift, section, and honest representation phase that these corrections were not authorized to change.

Each stage has a certificate and a failure mode. Combining them into one phrase such as ``equivariant attachment'' would hide which input is missing. In the equal-capacity case, even the first quadratic stage needs the native Arf match. In a source problem with a fixed embedding, the transvection may be disallowed. In even-order action problems, the 1-cocycle may survive. In all cases, the same arithmetic candidate must remain the subject of the construction.

\source{The source marking note leaves nontrivial stabilizer actions outside its established attachment theorem. This section adds a conditional finite group-action correction with a complete averaging proof and explicit even-order control. It is not a claim to have constructed general arithmetic stabilizers or solved the honest phase equations.}

\section{Audited claim map and the next arithmetic obligations}\label{sec:register}
\subsection{What this continuation adds to the progress paper}
The conductor branch now has a single normalization map, a proof of the necessary period divisibility, an expanded source self-pair calculation, and a precise conditional rank-nine filter. The actual-order certificate remains visible. The full and skipped rank-eight profiles show why even several correct conductor invariants cannot establish the needed native quotient.

The classical induction branch now identifies its field tower, tower-independence input, tame descent and central normalization separately. Its characteristic extension is named, and the failed fresh access to the underlying positive-characteristic source is not concealed. This leaves a meaningful selected input audit without overstating it as a new proof of all local correspondence results.

The marking branch now contains the full center-fixing extension criterion, the dihedral/quaternion distinction, the sharp free-embedding capacity test, and the exact rank-eight Arf condition. The explicit transvection and odd-order action corrections then explain how two later finite discrepancies can be repaired when their actual inputs and allowed changes are supplied.

\begin{longtable}{@{}>{\raggedright\arraybackslash}p{.23\textwidth}>{\raggedright\arraybackslash}p{.35\textwidth}>{\raggedright\arraybackslash}p{.36\textwidth}@{}}
\toprule
Statement & Status here & Boundary retained\\\midrule\endhead
Principal period and conductor & Ordinary derivation with [BH99] input & Does not construct a stratum\\[3pt]
Source self-pair $232$ & Expanded inherited ramification argument & Original source independence is retained\\[3pt]
Critical-four filter & Conditional actual-order consequence & Passing is not type occurrence or H5\\[3pt]
Rank-eight profile separation & Audited inherited construction and exact grade calculations & Three invariants do not force the native profile\\[3pt]
Classical source candidate & Selected primary-interface audit & Detailed [HL11] reread unavailable in this block\\[3pt]
Fixed-center extension criterion & Complete finite proof & Full sections can impose additional pointwise relations\\[3pt]
Free quadratic embedding & Complete finite proof & Requires the actual host refinement as input\\[3pt]
Transvection correction & Complete finite proof and exact tests & May change a preassigned embedding\\[3pt]
Odd-action correction & Complete conditional finite proof & Does not construct arithmetic actions or honest phases\\[3pt]
Native rank-eight refinement & Explicit remaining arithmetic obligation & Its value is not inferred from capacity\\[3pt]
General realization and H5 & Open in this project & No replacement by profile matching\\\bottomrule
\end{longtable}

\subsection{A certificate for the next native-marking step}
For a fixed arithmetic candidate at rank eight, a useful next certificate must identify the actual finite central extension or specified finite subgroup in its scalar pushout. It must give the chosen lift of each of the six symplectic basis vectors and compute its square under the full central character. These six bits determine the refinement once the commutator basis is fixed. Their three pairwise products determine the Arf sign through \eqref{eq:gauss}.

The certificate must also say whether the target marking is independently fixed or is allowed to be transported from the source. If it is fixed, compute the defect of the required quotient map. A vanishing defect directly certifies that map. An isotropic nonzero defect supplies the transvection only when changing the map is allowed. Every additional section element and transport lift then needs its actual multiplication check.

This task cannot be replaced by evaluating a canonical-looking quadratic form on an abstract six-dimensional space. That would choose new data rather than calculate the data of the specified arithmetic object. It also cannot be replaced by choosing another candidate with the same conductors. The candidate and full inducing character must remain fixed throughout.

\subsection{Verification and reproducibility boundaries}
The new standard-library exact suite contains 23 test methods. Its conductor tests cover the normalization table, grade dimensions, period divisibility, the rank-nine formal-degree identities, the rank-eight profiles and the finite field used for the projective source characters. Its marking tests cover extension group laws, square and section controls, Gauss signs, all labeled symplectic bases in dimension four, all linear-defect corrections through dimension six, relative and equivariant corrections, and odd versus even action cocycles.

The inherited 15-method Block 03 finite suite and 12-method Block 01 suite are replayed as separate regression gates. Their successful execution is not a new run of every historical arithmetic program. The supplied replay record lists exactly which commands were executed. Finite enumeration counts are not converted into a percentage of conjecture completion.

The source archive preserves the earlier compact checkpoint and versioned ledger history. It does not repair the missing Block 02 editable archive by inventing its contents. The v0.45 terminal compilation receipt remains unrecovered. No fresh Lean or Lake build occurred in this block.

The dedicated conjecture visual companion is the next publication task. Its mathematical route is now clear: normalized conductors, rejected and surviving rank-nine candidates, equal invariants with unequal capacity, transported versus native markings, the equal-capacity Arf gate, and the two conditional finite corrections. The original XI monograph, visual reader and Lean guide can then be synchronized against this stable claim map without presenting the general conjecture as a theorem.

\needspace{10\baselineskip}
\subsection{Source register}
The package's machine-readable source map supplies paths and hashes for inherited records. Page numbers below are printed pages unless a PDF-page locator is explicitly stated. Web-viewed primary sources are recorded as viewed sources, not as locally downloaded byte-pinned copies.

\begin{description}[style=nextline,leftmargin=1.5cm,labelwidth=1.35cm,font=\normalfont\bfseries,before=\raggedright]
\item[{[Q13]}] \emph{Adjoint Conductor Test161}. Retained dependency of Block166. The source invariant and conditional period-nine actual-order certificate. SHA-256 and exact source path are in the source map.
\item[{[Q14]}] \emph{Principal Order Conductor161}. Principal grades, leading-level normalization and necessary source matching.
\item[{[Q15]}] \emph{Arithmetic Constructor Audit}. The original source-defined induction and tame-descent branch, with full-character and source-input boundaries.
\item[{[Q16]}] \emph{Dyadic Markings167}. Full central extensions, transported finite markings, fixed native refinements and the sharp capacity criterion.
\item[{[D167]}] \emph{Dyadic Capacity167}. Full and skipped defining profiles, floor/ceiling grade counts, and the rank-eight shared-invariant example.
\item[{[BK]}] Colin J. Bushnell and Philip C. Kutzko. \emph{The Admissible Dual of GL(N) via Compact Open Subgroups}. Annals of Mathematics Studies 129, Princeton University Press, 1993. Relevant inherited inputs: recursive character orders, minimal commutator control, refinement, 6.2.5 and compact-type formal degree. The earlier selected audit remains distinguished from this block's fresh reading.
\item[{[BH99]}] Colin J. Bushnell and Guy Henniart. \emph{Local tame lifting for GL(n) II: wildly ramified supercuspidals}. Ast\'erisque 254, 1999. Freshly checked: introduction p. 2, 1.3 p. 15, 1.7 p. 18, 2.3 pp. 23--25, 6.1 pp. 53--54. The PDF page index is the printed page plus four when counted from zero.
\item[{[AP]}] Anne-Marie Aubert and Roger Plymen. \emph{Plancherel measure for GL(n,F) and GL(m,D): explicit formulas and Bernstein decomposition}. arXiv:math/0302169v4, 4 October 2004. Used: formulas (17), (29), Theorems 6.5 and 6.9, Remark 6.7 and the conductor convention in Section 4. The formal-degree pages inspected are pp. 19, 26--28.
\item[{[BH13]}] Colin J. Bushnell and Guy Henniart. \emph{To an effective local Langlands correspondence}. arXiv:1103.5316v2, 9 October 2013. The characteristic note on p. 8 is the selected input. Its effective construction is not used to supply its own totally wild starting correspondence.
\item[{[HL11]}] Guy Henniart and Bertrand Lemaire. \emph{Changement de base et induction automorphe pour GL(n) en caract\'eristique non nulle}. M\'emoires de la Soci\'et\'e Math\'ematique de France 124, 2011. DOI: 10.24033/msmf.436. The bibliographic record was checked. Detailed statements retain inherited audit status because the full PDF could not be retrieved in this block.
\end{description}

\end{document}
